% ch20: 泰勒公式与洛必达法则
\chapter{泰勒公式与洛必达法则}
\begin{applicationbox}
学完本章：用多项式逼近函数、用导数求极限
\end{applicationbox}

\section{洛必达法则}
若 \(\lim\frac{f(x)}{g(x)}\) 是 \(\frac00\) 或 \(\frac\infty\infty\) 型，且 \(\lim\frac{f'(x)}{g'(x)}\) 存在，则：
\[\lim\frac{f(x)}{g(x)} = \lim\frac{f'(x)}{g'(x)}\]

\textbf{ε-δ 证明（\(\frac00\) 型）：} 设 \(f(a)=g(a)=0\)，\(\lim_{x\to a}\frac{f'(x)}{g'(x)}=L\)。
对任意 \(\varepsilon>0\)，存在 \(\delta>0\) 使 \(0<|x-a|<\delta\) 时 \(|\frac{f'(x)}{g'(x)}-L|<\varepsilon\)。
对 \(x\in(a,a+\delta)\)，由柯西中值定理，存在 \(c\in(a,x)\) 使
\(\frac{f(x)-f(a)}{g(x)-g(a)}=\frac{f(x)}{g(x)}=\frac{f'(c)}{g'(c)}\)，
故 \(|\frac{f(x)}{g(x)}-L|=|\frac{f'(c)}{g'(c)}-L|<\varepsilon\)。\(\blacksquare\)

\begin{examplebox}[1]
\(\lim_{x\to0}\frac{\sin x}{x} = \lim_{x\to0}\frac{\cos x}{1} = 1\)
\end{examplebox}

\section{泰勒公式}
\(f(x) = f(a) + f'(a)(x-a) + \frac{f''(a)}{2!}(x-a)^2 + \cdots + \frac{f^{(n)}(a)}{n!}(x-a)^n + R_n(x)\)

\textbf{拉格朗日余项（ε-δ 证明框架）：} 令 \(F(t)=f(t)+\sum_{k=1}^{n-1}\frac{f^{(k)}(t)}{k!}(x-t)^k\)，
\(G(t)=(x-t)^n\)。由柯西中值定理，存在 \(\xi\) 在 \(a\) 与 \(x\) 之间使
\[R_n(x)=\frac{f^{(n)}(\xi)}{n!}(x-a)^n\]
即余项可表示为高阶导数在中间点的值乘以幂次——这是泰勒定理的精确形式。

\begin{examplebox}[2——常用展开]
\[e^x=1+x+\frac{x^2}{2!}+\frac{x^3}{3!}+\cdots,\quad \sin x = x-\frac{x^3}{3!}+\frac{x^5}{5!}-\cdots\]
\end{examplebox}

\section{习题}
\begin{exercise}[0] 洛必达法则适用于哪两种不定式？\end{exercise}
\begin{exercise}[0] \(\lim_{x\to0}\frac{\sin x}{x}\) 用洛必达怎么算？\end{exercise}
\begin{exercise}[0] \(e^x\) 在 \(x=0\) 处的泰勒展开第一项是什么？\end{exercise}
\begin{exercise}[0] \(\sin x\) 的泰勒展开的第一项是什么？\end{exercise}
\begin{exercise}[1] 求 \(\lim_{x\to0}\frac{e^x-1}{x}\)。\end{exercise}
\begin{exercise}[1] 求 \(\lim_{x\to0}\frac{\ln(1+x)}{x}\)。\end{exercise}
\begin{exercise}[1] 求 \(\lim_{x\to1}\frac{\ln x}{x-1}\)。\end{exercise}
\begin{exercise}[1] 写出 \(e^x\) 在 \(x=0\) 处的3次泰勒多项式。\end{exercise}
\begin{exercise}[2] 求 \(\lim_{x\to0}\frac{1-\cos x}{x^2}\)。\end{exercise}
\begin{exercise}[2] 求 \(\lim_{x\to\infty}\frac{x}{e^x}\)。\end{exercise}
\begin{exercise}[2] 求 \(\lim_{x\to0}\frac{\tan x}{x}\)。\end{exercise}
\begin{exercise}[2] 写出 \(\sin x\) 在 \(x=0\) 处的3次泰勒多项式。\end{exercise}
\begin{exercise}[3] 求 \(\lim_{x\to0}\frac{x-\sin x}{x^3}\)。\end{exercise}
\begin{exercise}[3] 求 \(\lim_{x\to0^+} x\ln x\)。\end{exercise}
\begin{exercise}[3] 求 \(\lim_{x\to\infty}\frac{\ln x}{x}\)。\end{exercise}
\begin{exercise}[3] 写出 \(\cos x\) 在 \(x=0\) 处的4次泰勒多项式。\end{exercise}
\begin{exercise}[3] 用泰勒展开求 \(e^{0.1}\) 的近似值（精确到0.001）。\end{exercise}
\begin{exercise}[4] 求 \(\lim_{x\to0}\frac{e^x-1-x}{x^2}\)。\end{exercise}
\begin{exercise}[4] 求 \(\lim_{x\to0}\frac{\arctan x - x}{x^3}\)。\end{exercise}
\begin{exercise}[4] 求 \(\lim_{x\to0}\left(\frac{1}{x} - \frac{1}{e^x-1}\right)\)。\end{exercise}
\begin{exercise}[4] 写出 \(\ln(1+x)\) 在 \(x=0\) 处的3次泰勒多项式。\end{exercise}
\begin{exercise}[4] 用二次泰勒展开近似 \(\sqrt{1.1}\)。\end{exercise}
\begin{exercise}[5] 求 \(\lim_{x\to0}\frac{\sin x - x + x^3/6}{x^5}\)。\end{exercise}
\begin{exercise}[5] 求 \(\lim_{x\to\infty}\left(1+\frac{1}{x}\right)^x\) 用洛必达法则。\end{exercise}
\begin{exercise}[5] 证明 \(\lim_{x\to0}\frac{x-\sin x}{x^3}=\frac16\)。\end{exercise}
\begin{exercise}[5] 用3次泰勒展开近似 \(\sin(0.1)\) 并估计误差。\end{exercise}
\begin{exercise}[6] 求 \(\lim_{x\to0^+} x^x\)。\end{exercise}
\begin{exercise}[6] 写出 \(\arctan x\) 在 \(x=0\) 处的5次泰勒展开。\end{exercise}
\begin{exercise}[6] 证明 \(e^x > 1 + x + x^2/2\)（\(x>0\)）。\end{exercise}
\begin{exercise}[7] 用泰勒展开证明 \(\sin x > x - x^3/6\)（\(x>0\)）。\end{exercise}
